\sectionheading{5.11 অনুবাদক প্রোগ্রাম (Translator Program)}

কম্পিউটার সরাসরি \textbf{Machine Language} বুঝতে পারে।
কিন্তু Programmer সাধারণত Assembly Language বা High Level
Language ব্যবহার করে Program লেখে।

এসব ভাষায় লেখা Program কম্পিউটার সরাসরি Execute করতে পারে না।
তাই Programmer-এর লেখা Program-কে কম্পিউটারের বোধগম্য
Machine Language-এ রূপান্তর করার জন্য একটি বিশেষ Program
প্রয়োজন হয়। এই Program-ই হলো \textbf{Translator Program}।

%================================================
% Definition
%================================================
\booksubheading{অনুবাদক প্রোগ্রাম কাকে বলে?}

\begin{definitionbox}

\textbf{\color{darkTeal}সংজ্ঞা}

যে Program কোনো Programming Language-এ লেখা Program-কে
কম্পিউটারের বোধগম্য \textbf{Machine Language বা Machine Code}-এ
অনুবাদ করে, তাকে
\textbf{অনুবাদক প্রোগ্রাম (Translator Program)} বলে।

\end{definitionbox}

\textbf{সহজভাবে মনে রাখবে—}

\begin{center}

\textbf{
Source Program
\(\longrightarrow\)
Translator Program
\(\longrightarrow\)
Machine/Object Code
}

\end{center}

%================================================
% Why Needed
%================================================
\booksubheading{Translator Program কেন প্রয়োজন?}

কম্পিউটারের Processor High Level Language বা Assembly Language-এর
Instruction সরাসরি বুঝতে পারে না।

তাই এসব ভাষায় লেখা Program-কে Machine-এর বোধগম্য রূপে পরিবর্তন
করতে Translator Program প্রয়োজন।

\begin{center}
\begin{adjustbox}{max width=0.98\linewidth}
\fcolorbox{darkBlue}{softBlue}{
    \parbox[c][9mm][c]{0.24\linewidth}{
        \centering\textbf{High Level /\\Assembly Code}
    }
}
\quad
\(\longrightarrow\)
\quad
\fcolorbox{mainTeal}{softGreen}{
    \parbox[c][9mm][c]{0.22\linewidth}{
        \centering\textbf{Translator\\Program}
    }
}
\quad
\(\longrightarrow\)
\quad
\fcolorbox{darkBlue}{softBlue}{
    \parbox[c][9mm][c]{0.23\linewidth}{
        \centering\textbf{Machine\\Code}
    }
}
\end{adjustbox}
\end{center}

অর্থাৎ Translator Program মানুষের জন্য সহজ Programming Language
এবং কম্পিউটারের Machine Language-এর মধ্যে
\textbf{যোগাযোগের মাধ্যম} হিসেবে কাজ করে।

%================================================
% Source Program
%================================================
\booksubheading{Source Program কী?}

\begin{definitionbox}

\textbf{\color{darkTeal}Source Program}

Programmer যে Programming Language ব্যবহার করে মূল Program
লেখে, তাকে \textbf{Source Program বা Source Code} বলা হয়।

\end{definitionbox}

\textbf{সহজভাবে—}

\begin{center}

\textbf{
Programmer-এর লেখা মূল Program
\(\;=\;\)
Source Program
}

\end{center}

যেমন Python, Java বা Assembly Language-এ Programmer যে Code
লেখে, সেটিই Source Program।

%================================================
% Object Program
%================================================
\booksubheading{Object Program কী?}

\begin{definitionbox}

\textbf{\color{darkTeal}Object Program}

Translator Program Source Program-কে অনুবাদ করার পর যে
Machine-oriented বা Target Program তৈরি করে, তাকে সাধারণভাবে
\textbf{Object Program বা Object Code} বলা হয়।

\end{definitionbox}

\begin{center}

\textbf{
Source Program
\(\longrightarrow\)
Translation
\(\longrightarrow\)
Object Program
}

\end{center}

\begin{rememberbox}

\textbf{\color{deepOrange}মনে রাখবে}

সব Translator একইভাবে আলাদা Object Program তৈরি করে না।

Compiler ও Assembler সাধারণত Object/Target Code তৈরি করে,
কিন্তু Interpreter-এর কাজের ধরন ভিন্ন।

\end{rememberbox}

%================================================
% General Flow
%================================================
\booksubheading{Translator Program-এর কাজের ধারা}

সাধারণভাবে Translator Program-এর কাজ—

\begin{center}
\begin{adjustbox}{max width=0.98\linewidth}
\fcolorbox{darkBlue}{softBlue}{
    \parbox[c][8mm][c]{0.19\linewidth}{
        \centering\textbf{Source\\Program}
    }
}
\(\longrightarrow\)
\fcolorbox{mainTeal}{softGreen}{
    \parbox[c][8mm][c]{0.20\linewidth}{
        \centering\textbf{Translator\\Program}
    }
}
\(\longrightarrow\)
\fcolorbox{darkBlue}{softBlue}{
    \parbox[c][8mm][c]{0.20\linewidth}{
        \centering\textbf{Machine /\\Object Code}
    }
}
\(\longrightarrow\)
\fcolorbox{mainTeal}{softGreen}{
    \parbox[c][8mm][c]{0.16\linewidth}{
        \centering\textbf{Execution}
    }
}
\end{adjustbox}
\end{center}

এখানে—

\begin{itemize}

\item \textbf{Source Program:}
Programmer-এর লেখা মূল Code।

\item \textbf{Translator Program:}
Source Code অনুবাদ করে।

\item \textbf{Machine/Object Code:}
কম্পিউটারের উপযোগী অনূদিত Code।

\item \textbf{Execution:}
কম্পিউটার Instruction অনুযায়ী কাজ সম্পন্ন করে।

\end{itemize}

%================================================
% Types
%================================================
\booksubheading{Translator Program-এর প্রধান প্রকারভেদ}

Translator Program প্রধানত তিন ধরনের—

\begin{tabularx}{\linewidth}{|>{\bfseries}p{0.22\linewidth}|Y|Y|}
\hline

\rowcolor{softBlue}
\multicolumn{1}{|c|}{\textbf{Translator}}
&
\multicolumn{1}{c|}{\textbf{কী অনুবাদ করে}}
&
\multicolumn{1}{c|}{\textbf{মূল ধারণা}}
\\
\hline

Compiler
&
High Level Language
&
সম্পূর্ণ Program একসঙ্গে অনুবাদ করে।
\\
\hline

Interpreter
&
High Level Language
&
সাধারণত Statement ধরে ধরে অনুবাদ ও Execute করে।
\\
\hline

Assembler
&
Assembly Language
&
Assembly Program-কে Machine Language-এ রূপান্তর করে।
\\
\hline

\end{tabularx}

%================================================
% Compiler Intro
%================================================
\booksubheading{Compiler-এর প্রাথমিক ধারণা}

Compiler High Level Language-এ লেখা
\textbf{সম্পূর্ণ Source Program একসঙ্গে} অনুবাদ করে।

\begin{center}

\textbf{
High Level Program
\(\longrightarrow\)
Compiler
\(\longrightarrow\)
Object/Target Program
}

\end{center}

Compiler পরবর্তী Topic-এ বিস্তারিত আলোচনা করা হবে।

%================================================
% Interpreter Intro
%================================================
\booksubheading{Interpreter-এর প্রাথমিক ধারণা}

Interpreter High Level Program-এর Instruction সাধারণত
\textbf{একটির পর একটি অনুবাদ করে এবং Execute করে}।

\begin{center}

\textbf{
Source Statement
\(\longrightarrow\)
Interpreter
\(\longrightarrow\)
Translation + Execution
}

\end{center}

Compiler ও Interpreter-এর পার্থক্য পরবর্তী অংশে বিস্তারিত
আলোচনা করা হবে।

%================================================
% Assembler Intro
%================================================
\booksubheading{Assembler-এর প্রাথমিক ধারণা}

Assembler হলো Assembly Language-এর Translator।

\begin{center}

\textbf{
Assembly Program
\(\longrightarrow\)
Assembler
\(\longrightarrow\)
Machine Language
}

\end{center}

\begin{rememberbox}

\textbf{\color{deepOrange}MCQ-এর জন্য গুরুত্বপূর্ণ}

\textbf{Assembler শুধু Assembly Language-এর Translator।}

\end{rememberbox}

%================================================
% Importance
%================================================
\booksubheading{Translator Program-এর প্রয়োজনীয়তা}

Translator Program-এর সাহায্যে—

\begin{compacttwocolpoints}

\item High Level Language-এ লেখা Program Machine-এর উপযোগী করা যায়।

\item Assembly Program Machine Code-এ রূপান্তর করা যায়।

\item Programmer-এর জন্য সহজ ভাষায় Program লেখা সম্ভব হয়।

\item Program-এর Error শনাক্ত করতে সহায়তা পাওয়া যায়।

\item Source Program ও Machine-এর মধ্যে প্রয়োজনীয় অনুবাদ সম্পন্ন হয়।

\end{compacttwocolpoints}

%================================================
% Source and Target Language
%================================================
\booksubheading{Source Language ও Target Language}

\begin{definitionbox}

\textbf{\color{darkTeal}Source Language}

যে ভাষায় মূল Program লেখা হয় এবং যেখান থেকে Program অনুবাদ করা হয়,
তাকে \textbf{Source Language} বলে।

\medskip

\textbf{\color{darkTeal}Target Language}

যে ভাষায় Program অনুবাদ করা হয়, তাকে
\textbf{Target Language} বলে।

\end{definitionbox}

উদাহরণ—

\begin{center}

\textbf{
Assembly Language
\(\longrightarrow\)
Assembler
\(\longrightarrow\)
Machine Language
}

\end{center}

এখানে—

\begin{itemize}

\item Assembly Language = \textbf{Source Language}

\item Machine Language = \textbf{Target Language}

\end{itemize}

%================================================
% Relationship
%================================================
\booksubheading{Translator Program ও Programmer-এর সম্পর্ক}

Programmer-কে সরাসরি 0 ও 1 ব্যবহার করে Machine Program লিখতে
হলে Programming অত্যন্ত কঠিন হতো।

Translator Program ব্যবহারের ফলে Programmer সহজ Programming
Language-এ Code লিখতে পারে এবং Translator সেই Code-কে
Machine-এর উপযোগী করে।

\begin{center}

\textbf{
সহজ ভাষায় Programming
\(\;+\;\)
Translator
\(\;=\;\)
Program বাস্তবায়ন সহজ
}

\end{center}

%================================================
% Comparison
%================================================
\booksubheading{Compiler, Interpreter ও Assembler এক নজরে}

\begin{tabularx}{\linewidth}{|>{\bfseries}p{0.18\linewidth}|Y|Y|Y|}
\hline

\rowcolor{softBlue}
\multicolumn{1}{|c|}{\textbf{বিষয়}}
&
\multicolumn{1}{c|}{\textbf{Compiler}}
&
\multicolumn{1}{c|}{\textbf{Interpreter}}
&
\multicolumn{1}{c|}{\textbf{Assembler}}
\\
\hline

Source Language
&
High Level
&
High Level
&
Assembly
\\
\hline

অনুবাদের ধরন
&
সম্পূর্ণ Program একসঙ্গে
&
সাধারণত Statement ধরে ধরে
&
Assembly থেকে Machine Code
\\
\hline

Object Code
&
সাধারণত তৈরি করে
&
সাধারণত পৃথক Object Program তৈরি করে না
&
তৈরি করে
\\
\hline

প্রধান কাজ
&
High Level থেকে Target Code
&
অনুবাদ ও Execute
&
Assembly থেকে Machine Code
\\
\hline

\end{tabularx}

%================================================
% Final Reminder
%================================================
\begin{rememberbox}

\textbf{\color{deepOrange}মনে রাখবে}

\begin{itemize}

\item Translator Program হলো \textbf{অনুবাদক প্রোগ্রাম}।

\item এর কাজ Source Program-কে Machine-এর উপযোগী Code-এ রূপান্তর করা।

\item Programmer-এর লেখা মূল Code হলো
\textbf{Source Program বা Source Code}।

\item অনুবাদের পর তৈরি Code-কে প্রসঙ্গভেদে
\textbf{Object Code বা Target Code} বলা হয়।

\item Translator Program প্রধানত তিন ধরনের—
\textbf{Compiler, Interpreter ও Assembler}।

\item Compiler ও Interpreter High Level Language-এর সঙ্গে সম্পর্কিত।

\item Assembler হলো \textbf{Assembly Language-এর Translator}।

\item
\textbf{
Assembly Language
\(\rightarrow\)
Assembler
\(\rightarrow\)
Machine Language
}

\item Compiler সম্পূর্ণ Program একসঙ্গে অনুবাদ করে।

\item Interpreter সাধারণত Statement ধরে ধরে অনুবাদ ও Execute করে।

\item Source Language হলো যে ভাষা থেকে অনুবাদ করা হয়।

\item Target Language হলো যে ভাষায় অনুবাদ করা হয়।

\end{itemize}

\end{rememberbox}

\small

%================================================
% 5.12 Compiler
%================================================
